\section{源调用的适配与统一表示}

四个源基准对任务、工具和参考调用的保存方式差异很大。$\tau$-bench 把任务写成 Python 构造器，BFCL 将多轮问题与参考答案分置于两组 JSONL，API-Bank 把一次对话拆成 User、AI 和 API 记录，AppWorld 则保存 HTTP 级方案执行日志。如果直接在这些文件上编写隐私构造逻辑，同一项检查会出现四套实现，来源语义也容易在格式转换中丢失。本文因此把数据构建分为两个阶段：Adapter 只负责恢复上游事实，公共 Builder 再负责档案、授权、筛选与配对调用生成。

\subsection{Adapter 的职责边界}

Adapter 是一个确定性的源格式解析器。它接收锁定的 benchmark 目录，输出一组内部任务对象及解析统计；其职责包括定位正式数据文件、恢复工具 schema、解析参考调用、确定当前任务可见的工具，并保存能够解释这些结果的来源路径。Adapter 不读取隐私字段配置，不计算 allowed/forbidden，不选择泄露步骤，也不根据调用数或敏感字段数筛选任务。由此，同一份 Adapter 输出可以在不重新解释上游数据的前提下供不同构造策略复用。

源版本检查位于 Adapter 之后、Builder 之前。系统根据 Adapter 声明的完整输入文件集合，核对源仓库 revision 与逐文件哈希。该顺序很重要：输入锁覆盖 Adapter 实际读取的文件，而不是预先假定的一组目录；新增或遗漏输入都会改变锁定结果。通过检查后，后续流程只接收已经绑定来源版本的任务对象。

\begin{figure}[t]
  \centering
  \begin{tikzpicture}[
    font=\small,
    box/.style={rounded corners=2pt, draw, thick, align=center, inner sep=5pt, minimum height=9mm},
    source/.style={box, draw=RPLBlue, fill=RPLBlue!5, text width=3.0cm},
    adapter/.style={box, draw=RPLGreen, fill=RPLGreen!6, text width=2.6cm},
    common/.style={box, draw=RPLOrange, fill=RPLOrange!7, text width=3.0cm},
    audit/.style={box, draw=gray, fill=gray!8, text width=3.0cm, font=\footnotesize},
    arrow/.style={-{Latex[length=2mm]}, thick},
  ]
    \node[source] (files) {锁定的源文件\\\footnotesize 任务·工具·状态·参考调用};
    \node[adapter, right=10mm of files] (adapter) {Source Adapter\\\footnotesize 解析与规范化};
    \node[common, right=10mm of adapter] (case) {\code{SourceCase}\\\footnotesize calls·tools·state·provenance};
    \node[common, right=10mm of case] (builder) {公共 Builder\\\footnotesize 筛选·授权·配对构造};
    \node[audit, below=9mm of adapter] (lock) {\code{SOURCE_LOCK.json}\\revision + input hashes};
    \node[audit, below=9mm of case] (contract) {Source contract\\origin + role + visibility};
    \draw[arrow] (files) -- (adapter);
    \draw[arrow] (adapter) -- (case);
    \draw[arrow] (case) -- (builder);
    \draw[arrow, dashed] (files) -- (lock);
    \draw[arrow, dashed] (adapter) -- (lock);
    \draw[arrow, dashed] (adapter) -- (contract);
    \draw[arrow, dashed] (contract) -- (case);
  \end{tikzpicture}
  \caption{Adapter 边界。上游文件先被解析为保留来源语义的 \texttt{SourceCase}，通过版本锁定后才进入公共构造流程；隐私规则不进入 Adapter。}
  \label{fig:adapter-boundary}
\end{figure}

\subsection{统一内部对象}

对任务 $c$，Adapter 输出

\begin{equation}
S_c=(d_c,x_c,C_c,T_c,U_c,M_c),
\end{equation}

其中 $d_c$ 是来源内的任务域，$x_c$ 是用户指令，$C_c=(a_1,\ldots,a_n)$ 是按来源顺序恢复的参考调用，$T_c$ 是该任务可见的工具集合，$U_c$ 是可用于恢复用户或环境状态的结构化数据，$M_c$ 保存文件、索引和来源角色等追溯信息。调用 $a_i=(t_i,g_i)$ 由规范工具名 $t_i$ 和参数映射 $g_i$ 组成。

代码中的 \code{SourceCase}、\code{SourceTrace}、\code{NormalizedCall} 和 \code{NormalizedTool} 分别承载这些信息。表\ref{tab:adapter-interface}列出公共 Builder 实际依赖的字段。这里的”统一”仅指访问接口一致。

\begin{table}[t]
  \centering
  \caption{Adapter 的统一输出接口及其保真要求。}
  \label{tab:adapter-interface}
  \small
  \begin{tabularx}{\textwidth}{L{3.1cm}L{3.2cm}Y}
    \toprule
    字段 & 内容 & 必须保持的性质 \\
    \midrule
    \code{source\_case\_id} & 上游任务标识或稳定派生标识 & 能唯一回到来源任务，不使用发布序号替代 \\
    \code{instruction} & 来源任务中的用户需求 & 仅做空白规范化，不补写任务意图 \\
    \code{trace.calls} & 规范工具名与参数映射 & 工具、参数值和调用顺序与来源记录一致 \\
    \code{trace.origin/role} & 调用来自何处、承担何种证据角色 & 区分目标动作、参考调用、聚合记录与执行日志 \\
    \code{turn\_boundaries} & 扁平调用序列中的轮次结束位置 & 有多轮结构时保留；来源未提供时写为 null \\
    \code{tools} & 当前任务可见的工具 schema & required、类型、返回值和可见范围可验证 \\
    \code{source\_state} & 用户档案或环境初态中的可用事实 & 不把无连接依据的局部身份合并为同一用户 \\
    \code{provenance} & 文件、split、行号、哈希或任务目录 & 足以复查每项适配决策 \\
    \bottomrule
  \end{tabularx}
\end{table}

统一对象在进入 Builder 前满足四项基本条件。第一，每个调用都能在 $T_c$ 中找到唯一 schema；第二，调用参数覆盖全部 required 参数并通过类型检查；第三，存在轮次边界时，边界单调不减且最后一个边界等于调用数；第四，\code{origin}、\code{role}、工具可见性依据和来源检查项均为非空的显式值。违反这些条件的任务在 Adapter 阶段拒绝，而不是带着不完整信息进入隐私构造。

\subsection{四类来源的规范化方法}

\subsubsection{$\tau$-bench：受限 AST 解析}

$\tau$-bench 的正式任务和工具定义是 Python 源码。Adapter 使用 \code{ast.parse} 检查 \code{Task(...)}、\code{Action(...)} 和工具 \code{get\_info()} 的语法树，只对预期节点调用字面量求值，不导入上游模块，也不执行其中的任意代码。每个 action 的名称必须属于当前领域工具表，参数必须符合工具声明。airline 与 retail 可能出现同名工具，因此规范名增加领域前缀，例如 \code{airline\_\_cancel\_reservation}。

任务去重键为 $(domain,user\_id,instruction)$。键相同且 actions 完全一致的记录折叠；同键但 actions 冲突的记录整体拒绝，避免由文件遍历顺序决定保留哪条轨迹。用户状态从对应领域的 \code{users.json} 读取。若参考调用显式指向唯一且不同的 \code{user\_id}，实体身份以调用目标为准，并在 provenance 中同时保存任务身份与解析后的实体身份；如果一条任务指向多个用户，则 Adapter 直接报错。

\subsubsection{BFCL：多轮边界与函数类作用域}

BFCL 以任务 ID 连接 question 和 possible answer。两侧 ID 集必须完全相同，重复 ID 或单侧缺失都会使转换失败。\code{possible\_answer.ground\_truth} 是按轮组织的函数调用字符串；Adapter 使用表达式模式的 AST 解析函数名、位置参数和关键字参数，不允许字典展开、未知参数或重复赋值。解析后的调用被扁平化，同时记录每一轮结束时的累计调用数，因此公共 Builder 可以顺序处理调用，又不会丢失原始轮次。

工具可见范围由 \code{involved\_classes} 决定。Adapter 只合并这些函数类的文档；若两个可见类对同名函数给出不同 schema，该名称被标记为歧义，不参与静默覆盖。\code{initial\_config} 原样进入 source state，供后续环境恢复和档案构造使用。它描述整个任务的初态，而不是某个单独轮次的输入。

\subsubsection{API-Bank：请求—结果绑定}

API-Bank 的工具 schema 从 \code{apis/} 下的 Python class 静态恢复。Adapter 读取类级 \code{description}、\code{input\_parameters} 和 \code{output\_parameters}，并根据 \code{call()} 函数签名区分必填参数。对话只接受角色为 User、AI 或 API 的 JSONL 行，其中只有显式 API 行进入参考调用。文件名和 AI 自然语言回复均不用于猜测工具调用。

每条 API 行必须同时给出 \code{param\_dict} 和 recorded \code{result}。安全类型规范化后，请求参数应与 \code{result.input} 的键和值一致，\code{result.exception} 必须为空，且参数仍需通过原生 schema。Level-2 还维护 ToolSearcher 已暴露工具集合：除 ToolSearcher 本身外，后续调用只有在之前的搜索结果中出现过才被接受。同一对话的 API 行按原顺序聚合为 $C_c$，User 行则确定 turn boundaries；原始请求和结果的摘要写入 provenance，便于检查聚合是否改变记录。

\subsubsection{AppWorld：HTTP 日志到工具调用}

AppWorld 的方案日志记录 HTTP method、URL 和 request data，而工具文档使用带占位符的 path template。Adapter 以 method 与路径模板进行唯一匹配，将 URL 中的路径参数按 schema 恢复为整数、数值、布尔值或字符串，再与 request data 合并。找不到路由、匹配多个路由或合并参数不符合 schema 时，任务均被拒绝。

当前任务可见工具取 \code{required\_apis.json} 与方案日志实际调用工具的并集。指令与任务用户来自 \code{specs.json}，private/public task data 进入 source state，并保留各自来源。\code{ground\_truth/api\_calls.json} 提供的是上游方案执行日志；Adapter 只恢复调用，不在转换阶段启动 AppWorld。独立环境执行及最终状态判断在评测阶段进行，避免把“日志成功解析”误写成“本项目已经执行”。

\subsection{SourceRules 与来源契约}

格式差异由 Adapter 处理，构造政策差异则集中在 \code{SourceRules}。每个来源注册一份契约，声明参考调用的 origin、role、可见工具依据、应出现的来源检查项以及是否包含轮次边界。Builder 在渲染报告时写入这些声明，Validator 再与注册值逐项比较，因此修改 Adapter 的解释口径会成为可检测的协议变化。

\code{SourceRules} 还提供三类来源专属函数：选择档案模板、从结构化状态提取字段，以及把工具归类为 internal、service provider、bank 或 public audience。与解析相关的事实留在 Adapter，与隐私构造相关的政策留在 SourceRules，二者没有通过大量条件分支混入公共 Builder。

\subsection{保守失败与 Reference 完整性}

适配阶段采用拒绝优先的策略。无法唯一确定工具、参数与 schema 不一致、来源身份冲突或请求—结果无法绑定时，不修写 reference，也不根据常识补齐调用。另有一类问题只能在原生环境中发现，例如工具返回“商品不存在”或业务状态不允许当前操作。这些任务仍可由 Adapter 忠实解析，但发布政策根据锁定的执行证据将其排除，并在 conversion report 中逐条记录 case ID 和原始错误。

发布链在参考序列前增加一次 \code{LoadEntityProfile}，因此链长为

\begin{equation}
L(c)=1+|C_c|.
\end{equation}

这里的 $C_c$ 始终是完整保留的来源序列。任务书中的 5/7 链形可以从 $|C_c|\in\{4,6\}$ 的任务派生，但不作为 Adapter 的截断规则。否则，格式要求会反过来改变来源证据，尤其会丢失 BFCL 的多轮调用和 AppWorld 的长执行日志。
